% Generated draft: requires mathematical and source review.
\documentclass[11pt]{article}
\usepackage[margin=25mm]{geometry}
\usepackage{fontspec}
\setmainfont{Latin Modern Roman}
\usepackage{amsmath,amssymb,mathtools,unicode-math}
\setmathfont{Latin Modern Math}
\usepackage{enumitem,needspace,xurl,hyperref}
\setlength{\parindent}{0pt}
\setlength{\parskip}{5pt}
\setlength{\emergencystretch}{3em}
\newcommand{\N}{\mathbb{N}}
\newcommand{\Z}{\mathbb{Z}}
\newcommand{\Q}{\mathbb{Q}}
\newcommand{\R}{\mathbb{R}}
\newcommand{\C}{\mathbb{C}}
\newcommand{\F}{\mathbb{F}}
\newcommand{\A}{\mathbb{A}}
\newcommand{\PP}{\mathbb{P}}
\newcommand{\E}{\mathbb{E}}
\newcommand{\eps}{\varepsilon}
\newcommand{\dd}{\,\mathrm{d}}
\begin{document}
\section*{UnsolvedMath ID: 20000378}
\textbf{Source problem code:} catalogue entry
\subsection{Definitions and mathematical statement}
Let $X\subset\PP^3_{\Q}$ be a smooth quartic K3 surface and $H$ a height for its hyperplane polarization. A geometrically rational curve means a geometrically integral curve whose normalization over $\overline\Q$ has genus zero. Being defined over $\Q$, having normalization $\PP^1_{\Q}$, and having a rational branch over a prescribed $x\in X(\Q)$ are different requirements. The source's phrase ``rational curve over $\Q$'' does not resolve them.

\textbf{Original Problem 10 and its attributed questions.} Peyre takes $U$ to be the complement of all rational curves over $\Q$, assumes $U(\Q)\ne\varnothing$, and reports numerical evidence for
\[\#\{u\in U(\Q):H(u)\le B\}\sim c(\log B)^{\rho(X)},\]
where $\rho(X)$ is the rank of the Picard group of $X$ and $c$ a product of local densities. Can such a scalar formula hold without weak approximation? This $U$ is a set-theoretic complement, not necessarily a Zariski open. The height, Picard-field convention and local-density normalization are not specified.

Tschinkel asks for a procedure deciding whether a given rational point of a quartic K3 lies on a rational curve over $\Q$. Skorobogatov asks about also removing elliptic curves. Browning asks for numerical evidence, referring to van Luijk, and asks for the Peyre freedom of small-degree rational curves on K3 surfaces. Colliot-Thélène asks whether a K3 surface over $\Q$ with a rational point has any rational curve over $\Q$, and whether one can find such a curve containing the prescribed point. Tschinkel asks the analogous existence and prescribed-point questions for $X/\F_q$, with curves defined over $\F_q$. The source reports that Bogomolov--Tschinkel obtain rational curves for Kummer surfaces over $\F_q$ for ``most'' points, without defining that qualification. Skorobogatov and Testa ask whether a K3 surface over $\Q$ can have infinitely many rational curves over $\Q$ while $\operatorname{Pic}(X_{\C})\cong\Z$.

\textbf{Attached unreviewed bounded-degree claim.} For an explicit smooth quartic $X/\Q$, a point $x\in X(\Q)$ and integer $D\ge1$, the aid claims a terminating Hilbert-scheme procedure listing the $\Q$-defined geometrically integral geometrically rational curves through $x$ of projective degree at most $D$. It separately decides descent of the Hilbert point, geometric genus zero, whether the genus-zero normalization is $\PP^1_{\Q}$, and whether that normalization has a rational point above $x$. It says a finite-field version detects Frobenius descent. A terminating answer for each fixed $D$ does not by itself decide the original unbounded-degree problem or termination of the search for a negative answer.

The aid further reports that a localized counting law with full support forces weak approximation. Precisely, the intended support condition requires positive leading constants for every nonempty product of open neighborhoods at a finite set of places; then such neighborhoods must contain rational points. A scalar total-count law alone gives no such distribution assertion. The precise measures and asymptotic uniformity are not supplied.

It also reports exact curve-degree and singularity constraints in the rank-one case $\operatorname{Pic}(X_{\C})=\Z H$, and a smooth quartic over $\F_5$ with no $\F_5$-defined geometrically rational curve. Neither the full constraints nor equations or verification for the finite-field example are supplied. All these results are descriptions of the machine-generated aid and require independent review.
\subsection{Short English statement}
The source combines quartic K3 counting, weak approximation, rational-curve decision and existence questions over $\Q$ and finite fields. The aid offers a bounded-degree procedure, carefully separating descent, geometric rationality, split normalization and a rational branch. Its localized full-support counting implication is stronger than the scalar asymptotic in the original question.
\subsection{Sources}
\begin{itemize}
\item Local UnsolvedMath registry, permanent ID 20000378, source code catalogue entry; \texttt{aim-algebraic-number-theory-notes.json}, notes[29], tag \texttt{problem}.
\item AIM workshop \emph{Rational and integral points on higher-dimensional varieties}, Problem 10: \url{https://aimath.org/pastworkshops/ratlhigherdimvarproblems.pdf}. Attributions: Peyre, Tschinkel, Skorobogatov, Browning, Colliot-Thélène, and Skorobogatov/Testa; references by name to van Luijk and Bogomolov--Tschinkel. No publication details are supplied.
\item The supplied verification says the record was checked against the workshop PDF and source TeX, but the replacement transcription is not included. This draft does not claim a new source visit. Original literature field empty.
\item Attached machine-generated research attempt 1: result \texttt{partial\_result}, status \texttt{partially\_solved}, rigor label \texttt{proved}, novelty confidence low. It is not peer reviewed and requires independent verification before citation. No locator for the procedure or finite-field example is supplied.
\end{itemize}
\end{document}
